\section{Parity subtraction from six privacy tests}\label{sec:general}\label{sec:construction}
The construction rests on a distinction between an exact algebraic
relation and the conditional entropy available for hashing. Four local bits
satisfying one parity relation have three algebraic degrees of freedom.
A common linear hash preserves that relation, while privacy amplification
preserves asymptotically one bit of conditional entropy per output bit
of each missing form. The following theorem gives a sufficient
condition for turning this structure into entropy subtraction.
All bit arithmetic below is over $\F$, with addition modulo two.

\begin{samepage}
\begin{theorem}[Parity-subtraction criterion]\label{thm:general}
Let $X_A,X_B,X_C,X_D$ be finite random variables with entropy vector $h$.
Suppose deterministic local labels $Y_i=\chi_i(X_i)\in\F$ satisfy
$Y_A+Y_B+Y_C+Y_D=0$ almost surely. Put
$U=Y_A$, $V=Y_B$, $E=Y_A+Y_C$, and let
\begin{align*}
 c=\min\{&H(E\mid X_AX_B),\ H(V\mid X_AX_C),\ H(V\mid X_AX_D),\\
         &H(U\mid X_BX_C),\ H(U\mid X_BX_D),\ H(U\mid X_CX_D)\}.
\end{align*}
Then, with $u(S)=\min\{3,|S|\}$,
\[
 h-Ru\in\AEcone\qquad(0\le R\le c).
\]
No uniformity or independence of the labels is assumed.
\end{theorem}
\end{samepage}

Each quantity defining $c$ is a privacy budget: how much randomness
remains in a missing bit after a pair's complete observations are
revealed. The minimum is the rate a single hash can achieve in all six
situations. In particular the criterion concerns the finite source and
its local labels, rather than its entropy vector alone.

\subsection{Algebraic completion and singleton privacy}
The parity identity writes the four labels as
\begin{equation}\label{eq:generalforms}
 (Y_A,Y_B,Y_C,Y_D)=(U,V,U+E,V+E).
\end{equation}
Their coefficient vectors in the formal basis $(U,V,E)$ are
$(1,0,0),(0,1,0),(1,0,1),(0,1,1)$.
Any two have rank two, and any three have rank three. These are statements
about the linear maps, even when the random inputs are dependent.
For every pair, Table~\ref{tab:basis} gives a missing form $L_{ij}$
and the explicit completion after hashing. Any other completing binary
linear form differs by a form known to the pair, so its conditional
entropy given $X_iX_j$ is the same.
\begin{table}[ht]
\centering
\begin{tabular}{@{}ccl@{}}
\toprule
Pair & $L_{ij}$ & $(K_A,K_B,K_C)$ from the pair's keys and $N_{ij}=ML_{ij}^n$\\
\midrule
$AB$ & $E$ & $(K_A,K_B,K_A+N_{ij})$\\
$AC$ & $V$ & $(K_A,N_{ij},K_C)$\\
$AD$ & $V$ & $(K_A,N_{ij},K_A+K_D+N_{ij})$\\
$BC$ & $U$ & $(N_{ij},K_B,K_C)$\\
$BD$ & $U$ & $(N_{ij},K_B,N_{ij}+K_B+K_D)$\\
$CD$ & $U$ & $(N_{ij},N_{ij}+K_C+K_D,K_C)$\\
\bottomrule
\end{tabular}
\caption{The six basis completions. Here $K_i=MY_i^n$ and
$K_D=K_A+K_B+K_C$. Conversely, $N_{AB}=K_A+K_C$,
$N_{AC}=N_{AD}=K_B$, and $N_{BC}=N_{BD}=N_{CD}=K_A$.
These are invertible changes of coordinates, with no distributional assumption.}
\label{tab:basis}
\end{table}

The six pair tests also guarantee the four singleton tests needed in the
proof. Conditioning reduces entropy, and~\eqref{eq:generalforms} gives
\begin{align}
 H(Y_B\mid X_A)&\ge H(V\mid X_AX_C)\ge c,\notag\\
 H(Y_A\mid X_B)&\ge H(U\mid X_BX_C)\ge c,\notag\\
 H(Y_D\mid X_C)&\ge H(Y_D\mid X_AX_C)=H(V\mid X_AX_C)\ge c,\label{eq:siximplyten}\\
 H(Y_C\mid X_D)&\ge H(Y_C\mid X_BX_D)=H(U\mid X_BX_D)\ge c.\notag
\end{align}
For the third line, $(X_A,X_C)$ fixes $E$, so $Y_D=V+E$ has the
same conditional entropy as $V$. The fourth line uses
$Y_C=U+Y_B+Y_D$, with $Y_B,Y_D$ fixed by $(X_B,X_D)$.
All side information here consists of the full observed variables.

\subsection{One linear hash for every test}
We use the following finite-alphabet form of conditional privacy
amplification~\cite{BBCM,Csiszar96,RW}. Its elementary proof is included
in Appendix~\ref{app:hash}.

\begin{samepage}
\begin{lemma}[Simultaneous linear hashing]\label{lem:hash}
Let $(Z_\ell,G_\ell)$, $1\le\ell\le q$, be a nonempty finite list of
functions on one finite
probability space, with $Z_\ell\in\F$ and $G_\ell$ finite. Suppose
$0\le R<\min_\ell H(Z_\ell\mid G_\ell)$.
Take $n$ independent copies of this joint space and let $k=\lfloor nR\rfloor$.
For a binary $k\times n$ matrix $M$ put
\[
 \Delta_\ell(M)=k-H(MZ_\ell^n\mid G_\ell^n)\ge0,
\]
where entropy is evaluated at the fixed matrix $M$.
There is a constant $C$, independent of $n$, such that for every $n$
some matrix $M_n$ satisfies $\sum_\ell\Delta_\ell(M_n)\le C$.
\end{lemma}
\end{samepage}

The pairs in the list may be arbitrarily dependent. A small normalised
deficit means that the hashed bit sequence retains nearly its full
entropy per source copy, even with the complete side information.
A random binary matrix makes two distinct strings collide with probability
$2^{-k}$. Jensen's inequality and this collision identity bound the
expected deficit; a conditional typicality estimate makes the bound
independent of $n$. Summing the nonnegative deficits selects one matrix
for the whole list. Using the same matrix is essential because unrelated
hashes would not preserve the parity identity. The proof in
Appendix~\ref{app:hash} follows the bounded-deficit route used in Lean.

\begin{lemma}[Conditional entropy profiles]\label{lem:cone}
The almost-entropic region is a closed convex cone. If
$X_A,X_B,X_C,X_D,W$ are finite random variables, then
$S\mapsto H(X_S\mid W)$ belongs to $\AEcone$.
\end{lemma}
\begin{proof}
The cone property is classical~\cite[Theorem~1]{ZY}; we recall the
argument needed here. Independent products add entropy vectors.
For an entropic vector $f$ and $\alpha>0$, take an integer $m>\alpha$,
$m$ independent realisations of $f$, and an independent flag with
probability $\alpha/m$ of being one. At each coordinate retain the flag
and, when it is one, all $m$ samples; otherwise retain a constant.
Each nonempty subset entropy is $\hb(\alpha/m)+\alpha f(S)$, which tends
to $\alpha f(S)$. Approximation and closedness extend addition and scaling
to the closure, giving convexity.
Finally
\[
 H(X_S\mid W)=\sum_{w:\Pr(W=w)>0}\Pr(W=w)H(X_S\mid W=w)
\]
is a finite convex combination of entropy vectors of conditional laws.
This also explains the contraction operation of~\cite[Lemma~1]{MC}.
\end{proof}

\subsection{Accounting for all fifteen entropy losses}
\begin{proof}[Proof of Theorem~\ref{thm:general}]
Null atoms may be discarded. If $c=0$ only $R=0$ is at issue; otherwise
fix $0<R<c$ and $k=\lfloor nR\rfloor$.
Form every block variable from the same $n$ independent copies of the
original four-variable source.
Apply Lemma~\ref{lem:hash} to the six pairs $(L_{ij},X_iX_j)$ and the
four label--singleton pairs in~\eqref{eq:siximplyten}. Select one common
matrix $M=M_n$. Define
\[
 K_i=MY_i^n,\qquad W=(K_A,K_B,K_C),
\]
and the nonnegative deficits
\begin{align}
 \Delta_{ij}&=k-H(ML_{ij}^{\,n}\mid X_i^nX_j^n),\notag\\
 \delta_i&=k-H(K_{j(i)}\mid X_i^n),\qquad
 j(A)=B,\ j(B)=A,\ j(C)=D,\ j(D)=C,\label{eq:deficits}\\
 D_n&=\sum_{ij}\Delta_{ij}+\sum_i\delta_i\le C.\notag
\end{align}
Here $ij$ ranges over unordered pairs.
The common map gives $K_D=K_A+K_B+K_C$ exactly.
For every pair $ij$, basis completion is an invertible change of
coordinates between $W$ and $(K_i,K_j,ML_{ij}^{\,n})$.
Since $K_i,K_j$ are known from $X_i^nX_j^n$, this gives the exact identity
\begin{equation}\label{eq:pairs}
 H(W\mid X_i^nX_j^n)=H(ML_{ij}^{\,n}\mid X_i^nX_j^n)=k-\Delta_{ij}.
\end{equation}

For a singleton $i$, use its partner $j=j(i)$ in~\eqref{eq:deficits}.
The same basis change and the chain rule give
\begin{align*}
 H(W\mid X_i^n)
 &=H(K_j\mid X_i^n)+H(ML_{ij}^{\,n}\mid X_i^n,K_j)\\
 &\ge k-\delta_i+H(ML_{ij}^{\,n}\mid X_i^n,X_j^n)
 =2k-\delta_i-\Delta_{ij}.
\end{align*}
The inequality holds because $K_j$ is a function of $X_j^n$.
The upper bound is $2k$, since $K_i$ is already known.
Any triple determines $W$, so larger observed sets leave no uncertainty.

To bound $H(W)$ itself, use the invertible change
$W\leftrightarrow(K_A,K_B,ME^n)$ and condition each chain-rule term
on the corresponding full observations:
\begin{align*}
 H(W)&=H(K_A)+H(K_B\mid K_A)+H(ME^n\mid K_A,K_B)\\
 &\ge H(K_A\mid X_B^n)+H(K_B\mid X_A^n)
                  +H(ME^n\mid X_A^nX_B^n)\\
 &=3k-\delta_B-\delta_A-\Delta_{AB}.
\end{align*}
Also $H(W)\le3k$. Thus $H(W)=3k-d_0$ and
$H(W\mid X_i^n)=2k-d_i$, with $0\le d_0,d_i\le D_n$.
The complete entropy accounting is
\begin{center}
\begin{tabular}{@{}lccc@{}}
\toprule
Observed subset size & $1$ & $2$ & $3$ or $4$\\
\midrule
$H(W\mid X_S^n)$ & $2k-d_i$ & $k-\Delta_{ij}$ & $0$\\
$H(W)-H(W\mid X_S^n)$ & $k+O(1)$ & $2k+O(1)$ & $3k+O(1)$\\
\bottomrule
\end{tabular}
\end{center}
Since $H(X_S^n)=nh(S)$ by independence of the copies, the identity
\[
 H(X_S^n\mid W)=nh(S)-H(W)+H(W\mid X_S^n)
\]
now proves, for all fifteen nonempty subsets,
\begin{equation}\label{eq:all15}
 \left|H(X_S^n\mid W)-\bigl(nh(S)-ku(S)\bigr)\right|\le D_n.
\end{equation}
The errors are $d_0-d_i$ for singletons, $d_0-\Delta_{ij}$ for pairs,
and $d_0$ for larger sets; each absolute value is at most $D_n$.

\smallskip\noindent\emph{Passing to an almost-entropic vector.}
For each finite $n$, let $p_w=\Pr(W=w)$ and let $h_w$ be the entropy
vector of the four blocks conditioned on $W=w$. The conditional profile
is the finite average
\[
 \bar h_n(S)=H(X_S^n\mid W)=\sum_{w:p_w>0}p_w h_w(S).
\]
Lemma~\ref{lem:cone} gives $\bar h_n/n\in\AEcone$. Dividing~\eqref{eq:all15} by $n$, with $k/n\to R$ and $D_n/n\to0$,
gives $\bar h_n/n\to h-Ru$ coordinatewise. Closedness proves membership
for $0<R<c$, then for $R=c$ by taking $R\uparrow c$. The case $R=0$
is immediate. This passage uses an average of conditional profiles;
a single conditioning value is not required to realise that average.
\end{proof}

